\subsection{학습 인프라}
\label{sec:train_infra}

24\,B 사용자 이벤트에 걸친 207\,B 토큰으로 PRAGMA를 사전학습하는 일은 몇 가지 엔지니어링 과제를 동반한다.
데이터의 이질적이고 테이블 구조적인 성격은, 그에 맞춘 저장·배치·절단 전략을 요구한다.
저자들은 아래에서 각 항목을 차례로 기술한다.

\paragraph*{데이터 저장.}
사전학습 코퍼스는 2단계 구조로 저장된다 — \emph{사용자 인덱스(user index)}(각 사용자를 그 사용자의 토큰화된 프로필 상태와 사용자별 토큰 통계로 매핑하는 LMDB 기반 키-값(key-value) 저장소)와, \emph{이벤트 샤드(shard)} 모음(이벤트 개수로 파티셔닝(partitioning)된 Parquet 파일들로, 각 파일에는 동일한 이벤트 개수를 가진 사용자만 들어 있다).
이 레이아웃은 워커가 이벤트 샤드를 독립적으로 스트리밍하면서 프로필 상태는 필요할 때만 조회하도록 해 준다.

\paragraph*{배치.}
각 학습 샘플은 완전한 이벤트 이력과 그에 대응되는 프로필 상태 토큰으로 구성된다.
이벤트 이력 길이는 몇 건에서 수천 건까지 크게 변동하므로, 패딩 기반 배치를 단순히 적용하면 대부분의 연산이 패딩 토큰에 낭비될 것이다.
이벤트 개수로 레코드를 샤딩하면 로딩 단계에서 디스크 임의접근이 줄고, 한 배치 안의 이벤트 시퀀스 길이가 균일해진다. 결과적으로 히스토리 인코더는 들쭉날쭉하거나 패딩된 차원이 없는 직사각 텐서 위에서 작동한다.
저자들은 GPU 메모리에 들어맞는 고정 토큰 예산을 갖는 \emph{동적 배치(dynamic batching)} 를 사용한다 — 같은 샤드의 레코드를 예산이 찰 때까지 그리디하게 채워 넣는 방식이다.

\paragraph*{시퀀스 패킹.}
한 배치 안에서도 개별 이벤트의 토큰 수는 여전히 다르다.
저자들은 가장 긴 이벤트에 맞춰 모든 이벤트를 패딩(padding)하는 대신, 모든 이벤트 토큰을 평탄(flat) 버퍼에 패킹하고 가변 길이(varlen) 어텐션 커널(kernel)~\citep{dao2022flashattention}로 처리해, 이 단계에서는 서로 다른 이벤트의 토큰이 서로 어텐션하지 않도록 한다.
샤드 기반 배치와 결합하면, 이는 이벤트 축과 토큰 축 양쪽의 패딩 오버헤드를 모두 제거한다.
패딩된 베이스라인과 비교했을 때, 시퀀스 패킹과 동적 배치를 함께 쓰면 데이터셋의 시퀀스 길이 분포에 따라 $2$--$5{\times}$의 처리량 향상이 나타난다.

\paragraph*{절단.}
고정 컨텍스트 길이에서 메모리 소비를 제한하기 위해, 저자들은 패킹 전에 두 단계의 절단을 적용한다.
\emph{이벤트 단계}에서는 각 개별 이벤트를 최대 24개 토큰으로 자르며, 이는 전체 이벤트의 0.01\,\%에만 영향을 준다.
\emph{프로필 상태 단계}에서는 정적 프로필 상태 시퀀스를 최대 200개 토큰으로 자른다.
이벤트가 0인 사용자는 폐기하고, 6{,}500개 이벤트를 초과하는 사용자는 가장 최근 이벤트들을 남기는 방식으로 서브샘플링하여 시간 최근성(temporal recency)을 보존한다.

\paragraph*{사전학습 컴퓨트.}
세 모델 변형은 모두 bf16 혼합정밀도(mixed precision)와, Muon 옵티마이저(optimiser)를 AdamW와 결합한 설정으로 학습되었다~\citep{loshchilov2019adamw,jordan2024muon,liu2025muon}.
PRAGMA-S(10\,M 파라미터)와 PRAGMA-M(100\,M)은 $16{\times}$\,NVIDIA H100 GPU에서, PRAGMA-L(1\,B)은 $32{\times}$\,NVIDIA H100 GPU에서 학습되었다.
가장 작은 변형은 약 2~일 만에 수렴한 반면, 100\,M와 1\,B 모델은 각각 약 2~주의 벽시계 시간을 필요로 했다.
